\chapter{Conclusiones} \label{cap:conclusiones}

\section{Resumen del Proyecto}
El desarrollo de GuardiApp ha permitido transformar un proceso crítico y tradicionalmente manual, como es la gestión de guardias médicas, en una solución digital moderna, segura y eficiente. Se ha logrado centralizar la información, facilitar la comunicación entre profesionales y automatizar tareas administrativas que consumían gran cantidad de tiempo.

\section{Objetivos Alcanzados}
Se han cumplido satisfactoriamente los objetivos planteados al inicio del proyecto:
\begin{itemize}
    \item \textbf{Digitalización total}: Eliminación del uso de papel y hojas de cálculo locales para el cuadrante.
    \item \textbf{Gestión de Intercambios}: Implementación de un flujo transparente y trazable para los cambios de turno.
    \item \textbf{Control de Presencia}: Validación efectiva de la jornada laboral mediante geolocalización.
    \item \textbf{Integración Tecnológica}: Uso exitoso de un stack moderno (Laravel + React) que garantiza un rendimiento óptimo.
\end{itemize}

\section{Valoración Personal}
Este proyecto ha supuesto un reto técnico significativo, especialmente en la integración de la API REST con una interfaz dinámica en React. La implementación de la lógica de asignación automática de jefes de guardia y el sistema de notificaciones en tiempo real han permitido profundizar en el desarrollo de aplicaciones web de alta complejidad. 

En conclusión, GuardiApp se presenta como una herramienta lista para ser utilizada en un entorno hospitalario real, mejorando la calidad de vida laboral de los profesionales sanitarios.
